LLM repair loops improve code generation quality by feeding execution output
back to the model for iterative correction --- but whether adding static analysis
feedback provides meaningful improvement over execution-only feedback has never
been tested in a controlled experiment. We conduct a five-sub-hypothesis study
on HumanEval+ (164 problems, GPT-4o-mini, k=5 repair rounds) to isolate the
marginal contribution of mypy type checking over execution-only feedback. Our
primary finding: execution+mypy repair achieves 87.20\% pass@1 versus 78.25\%
for execution-only --- a consistent +8.94 percentage point improvement across
three independent seeds. Strikingly, this improvement is \emph{not} type-specific:
execution-only repair achieves identical 90.9\% repair rates on type-error
problems (differential = 0.000), consistent with mypy functioning as a general
diagnostic context enricher rather than a type-targeted correction signal
(this mechanism remains unconfirmed via ablation; see Section~6).
We further show that HumanEval+ and MBPP+ have fundamentally different mypy
error profiles (70\% vs 0\% in failing solutions), predicting where static
analysis feedback will and will not activate. A Z3 formal verification
extension achieves 91.1\% spec validity but provides no additional pass@1
improvement at the current counterexample activation rate (16.1\%). Our
results offer the first controlled evidence that one structured feedback
channel --- a single mypy subprocess call --- substantially improves LLM repair
quality, while revealing that the intuitive type-specificity mechanism is not
the driver.
